\documentclass[11pt,a4paper]{article}
\usepackage[utf8]{inputenc}
\usepackage[T1]{fontenc}
%\usepackage[italian]{babel}
\usepackage[margin=2.6cm]{geometry}
%\usepackage{lmodern}
%\usepackage{microtype}
\usepackage{booktabs}
\usepackage{amsmath}
\usepackage{amssymb}
\usepackage{array}
\usepackage{xcolor}
\usepackage{listings}
\usepackage{hyperref}
\usepackage{enumitem}
\hypersetup{colorlinks=true, linkcolor=blue!50!black, urlcolor=blue!50!black}

\lstset{
  basicstyle=\ttfamily\small,
  breaklines=true,
  frame=single,
  framerule=0.2pt,
  backgroundcolor=\color{gray!6},
  keywordstyle=\color{blue!60!black},
  commentstyle=\color{green!40!black},
  showstringspaces=false,
  language=Python
}

\newcommand{\code}[1]{\texttt{#1}}
\newcommand{\flow}[1]{\texttt{\small #1}}

\title{\textbf{Pipeline SDMX--ISTAT per popolazioni sintetiche comunali}\\[2mm]
\large Note tecniche --- v0.1}
\author{Mirko~Degli Esposti}
\date{18 luglio 2026}

\begin{document}
\maketitle

\begin{abstract}
\noindent
Queste note documentano la costruzione e la validazione di una pipeline automatica
per l'estrazione di vincoli demografici e socio-economici a livello comunale dal
web service SDMX di ISTAT (\code{esploradati.istat.it}), finalizzata alla sintesi
di popolazioni MaxEnt (GibbsPCD). La pipeline --- quattro script Python eseguiti
in locale (WSL2, env \code{ml}) --- copre: catalogazione e ricerca sistematica dei
4\,884 dataflow disponibili; fetch generico guidato dal DSD; costruzione del
constraint set per un comune arbitrario. Validazione completa sul caso Brescia
(017029, ancoraggio 1/1/2025). Risultato metodologico principale: a livello
comunale, anagrafe post-censuaria e censimento permanente coincidono
\emph{esattamente} cella per cella su sesso$\times$et\`a (MAE $=0$), e le stime
censuarie via SDMX sono servite non arrotondate; la distinzione hard/soft dei
vincoli va quindi ridisegnata.
\end{abstract}

\tableofcontents

%=======================================================================
\section{Scopo e contesto}
Obiettivo: una pipeline riproducibile che, dato un codice comune ISTAT, produca il
miglior constraint set ottenibile dai dati pubblici per la sintesi di una
popolazione individuale (attributi: sesso, et\`a a singolo anno, stato civile,
cittadinanza, istruzione, condizione professionale, nazionalit\`a). Il caso guida
\`e Brescia (SIN Caffaro), con estensione prevista all'attributo di zona
sub-comunale (33 quartieri). Il lavoro sostituisce il precedente flusso su Colab
con esecuzione locale (WSL2 Ubuntu, conda env \code{ml}, storage nativo in
\code{\textasciitilde/progetti/}).

%=======================================================================
\section{Endpoint SDMX \code{esploradati}: comportamenti non documentati}
\label{sec:quirks}
Lezioni operative emerse, tutte incorporate come guardie nel modulo
\code{istat\_sdmx.py}:

\begin{enumerate}[leftmargin=*]
  \item \textbf{Limite di 260 caratteri per segmento URL.} Il frontend Microsoft
  (http.sys) rifiuta con \code{400 Bad Request -- Invalid URL} qualunque richiesta
  il cui KEY superi $\sim$260 caratteri come singolo segmento di path ---
  \emph{prima} che la richiesta raggiunga l'applicazione SDMX. Il query builder
  del sito genera URL che violano questo limite (es.\ lista esplicita di 101
  et\`a). Soluzione: wildcard (posizione vuota nel KEY) e filtro dei codici
  aggregati a valle.
  \item \textbf{\code{startPeriod}/\code{endPeriod} ignorati} su (almeno alcune)
  dataflow annuali: il server restituisce l'intera serie storica. Il filtro
  temporale va fatto client-side su \code{TIME\_PERIOD}. Effetto collaterale
  positivo: serie storiche complete gratis (anagrafe 2019--2025, censimento
  2018--2024).
  \item \textbf{SDMX-CSV via content negotiation.} Con header
  \code{Accept: application/vnd.sdmx.data+csv;version=1.0.0} il server serve CSV
  direttamente ingeribile da pandas; niente parsing SDMX-ML per i dati.
  \item \textbf{Rate limit} $\sim$5 query/minuto: throttling globale a 13\,s tra
  richieste; timeout sporadici (retry con backoff, 3 tentativi).
  \item \textbf{Stime censuarie non arrotondate.} A differenza delle tavole web
  (che dichiarano arrotondamenti), il canale SDMX serve le stime del censimento
  permanente con i decimali, internamente additive: i residui
  totale$-$somma delle modalit\`a sono \emph{esattamente zero}
  (cfr.~\S\ref{sec:key-result}).
\end{enumerate}

%=======================================================================
\section{I tre ``laghi'' del datalake ISTAT}
\label{sec:laghi}
La reperibilit\`a degli attributi a livello comunale \`e determinata dalla
struttura della disseminazione, articolata in tre fonti con propriet\`a diverse:

\begin{center}
\begin{tabular}{@{}p{3.1cm}p{4.6cm}p{2.6cm}p{3.6cm}@{}}
\toprule
\textbf{Fonte} & \textbf{Contenuto comunale} & \textbf{Et\`a} & \textbf{Natura}\\
\midrule
Anagrafe / ANPR (\flow{DCIS\_*}) & sesso, et\`a, stato civile; stranieri; paese di cittadinanza & singolo anno & conteggi amministrativi, riconciliati col censimento\\
\addlinespace
Censimento permanente (\flow{DF\_DCSS\_*}) & istruzione, condizione professionale, cittadinanza, famiglie, background migratorio, pendolarismo & singolo anno o classi, secondo tavola & stime (campione + fonti amministrative), annuali dal 2018\\
\addlinespace
Indagini campionarie (AVQ, EHIS, \dots) & --- (solo regionale/provinciale) & classi & stime campionarie\\
\bottomrule
\end{tabular}
\end{center}

\noindent
Conseguenza architetturale: \emph{spine} anagrafica + \emph{strato} censuario a
livello comune; salute e stili di vita entrano solo come condizionali regionali
(approccio two-stage, coerente con quanto gi\`a fatto per la struttura geografica
in Brescia v3).

%=======================================================================
\section{Catalogo dataflow: la ricerca sistematica}
Il problema del ``non trovare'' i dati si risolve scaricando una volta l'elenco
completo dei dataflow (\code{/dataflow/IT1}: 4\,884 voci, 13\,MB XML) e
greppandolo in locale (\code{istat\_catalog.py}). Osservazioni:
\begin{itemize}[leftmargin=*]
  \item le tavole censuarie \emph{non} contengono ``censimento'' nel nome: si
  riconoscono dal prefisso \flow{DF\_DCSS\_*};
  \item la granularit\`a \`e codificata nel suffisso/nome: \flow{- comuni}
  (tutti i comuni), \flow{\_GC} (province e grandi comuni), \flow{comuni con pi\`u
  di 20.000 ab.};
  \item il grep per parola chiave (\code{istruzione}, \code{condizione prof},
  \code{cittadin}, \code{stranier}) individua in pochi secondi la rosa dei
  dataflow utili.
\end{itemize}

%=======================================================================
\section{Architettura della pipeline}
Quattro script in \code{\textasciitilde/progetti/gsp/scripts/}:

\begin{center}
\begin{tabular}{@{}lp{9.8cm}@{}}
\toprule
\textbf{Script} & \textbf{Funzione}\\
\midrule
\code{istat\_catalog.py} & scarica/cachea il catalogo dataflow; \code{grep\_catalog(pattern)}\\
\code{istat\_sdmx.py} & modulo generico: cache strutture per dataflow, dimensioni dal DSD, \code{build\_key} con wildcard, fetch SDMX-CSV, decodifica etichette via mapping dimensione$\to$codelist del DSD, throttling e retry\\
\code{fetch\_comune.py} & driver: rosa dei fondamentali per un codice comune; individuazione automatica della dimensione territorio; modalit\`a \code{--explore} (solo strutture) / fetch / \code{--only}; cache dei CSV scaricati\\
\code{build\_constraints\_v1.py} & dai CSV decodificati al constraint set (\S\ref{sec:constraints}); parametrico su comune e anno\\
\bottomrule
\end{tabular}
\end{center}

\noindent Scelte di design rilevanti:
\begin{itemize}[leftmargin=*]
  \item \textbf{KEY guidato dal DSD}: l'ordine delle dimensioni \`e letto dalla
  struttura, mai hardcodato; le dimensioni non specificate vanno in wildcard.
  \item \textbf{Decodifica via DSD, non euristica}: il mapping
  dimensione$\to$codelist \`e quello dichiarato (una prima euristica per
  inclusione di codici agganciava \flow{CL\_ITTER107} alla dimensione sesso).
  \item \textbf{Territorio auto-rilevato}: per ogni tavola si cerca la dimensione
  la cui codelist contiene il codice comune; robusto a nomi diversi
  (\code{REF\_AREA}, \code{ITTER107}, \dots).
  \item \textbf{Esecuzione staged}: explore (strutture) $\to$ smoke test (una
  tavola) $\to$ run completo; ogni fase validata prima della successiva.
\end{itemize}

%=======================================================================
\section{L'ipercubo DCSS}
\label{sec:ipercubo}
Tutte le tavole censuarie comunali condividono un unico DSD con $\sim$10
dimensioni (\code{INDICATOR}, \code{GENDER}, \code{AGE\_NOCLASS},
\code{CITIZENSHIP}, \code{EDU\_ATTAIN}, \code{CUR\_ACT\_STAT}, \code{LOC\_DEST},
\code{REAS\_COMMUTING}, \dots): ogni ``tavola'' \`e una fetta dell'ipercubo, con
le dimensioni non pertinenti fissate al codice totale
(\code{TOTAL}/\code{ALL}/\code{99}). Implicazioni:
\begin{itemize}[leftmargin=*]
  \item le cardinalit\`a delle codelist (es.\ 343 codici et\`a, inclusi i giorni
  di vita) sono della codelist condivisa, non della tavola; i valori effettivi li
  determina il ContentConstraint;
  \item i CSV contengono molte colonne costanti al totale: firma dell'ipercubo,
  da collassare nel builder;
  \item il wildcard su un singolo comune produce risposte piccole
  (10$^2$--10$^3$ righe/anno).
\end{itemize}

%=======================================================================
\section{Rosa dei fondamentali e semantica dei codici}
\label{sec:rosa}
Tavole scaricate per comune (righe riferite a Brescia 017029, tutte le
annualit\`a):

\begin{center}
\small
\begin{tabular}{@{}llllr@{}}
\toprule
\textbf{Nome} & \textbf{Dataflow} & \textbf{Incrocio} & \textbf{Et\`a} & \textbf{Righe}\\
\midrule
anag\_sesso\_eta\_statociv & \flow{22\_289\_DF\_DCIS\_POPRES1\_26} & S$\times$E$\times$SC & singola & 8\,568\\
cens\_sesso\_eta\_cittadinanza & \flow{DF\_DCSS\_POP\_DEMCITMIG\_SETA\_1} & S$\times$E$\times$C & singola & 3\,658\\
cens\_istruzione\_eta & \flow{DF\_DCSS\_ISTR\_LAV\_PEN\_2\_TV\_1} & S$\times$E$\times$I & 4 classi & 819\\
cens\_istruzione\_cittadinanza & \flow{DF\_DCSS\_ISTR\_LAV\_PEN\_2\_TV\_2} & C$\times$I & tot.\ (9+) & 441\\
cens\_condprof\_eta & \flow{DF\_DCSS\_ISTR\_LAV\_PEN\_2\_TV\_3} & S$\times$E$\times$P & 4 classi & 819\\
cens\_condprof\_cittadinanza & \flow{DF\_DCSS\_ISTR\_LAV\_PEN\_2\_TV\_4} & C$\times$P & tot.\ (15+) & 486\\
cens\_stranieri\_paesi & \flow{DF\_DCSS\_POP\_DEMCITMIG\_TV\_3} & paese & --- & 2\,454\\
\bottomrule
\end{tabular}
\end{center}

\noindent Semantica dei codici verificata sul dato (non ovvia dalle codelist):
\begin{itemize}[leftmargin=*]
  \item \textbf{Cittadinanza}: il codice effettivo per gli stranieri \`e
  \code{FRGAPO} (stranieri/apolidi), non \code{FRG}; normalizzato a \code{FRG} in
  uscita. Fallback implementato: derivazione per differenza
  \code{TOTAL}$-$\code{ITL}.
  \item \textbf{Condizione professionale, gerarchia a due livelli} (verificata
  additiva al centesimo): $22 = 1 + 12$ (forze di lavoro = occupati + in cerca);
  $23 = 4+5+7+24$ (non forze di lavoro = casalinga/o + studente/ssa + altra
  condizione + percettore di pensioni/redditi da capitale). Gli aggregati 22/23
  vengono scartati \emph{solo se} il rispettivo dettaglio \`e presente nel dato.
  \item \textbf{Classi d'et\`a per regola di contenimento}: un codice classe \`e
  aggregato se il suo intervallo contiene interamente un'altra classe osservata
  (cos\`i \code{Y\_GE9} \`e aggregato, \code{Y\_GE65} \`e elementare); nessun
  hardcoding per comune o tavola.
\end{itemize}

%=======================================================================
\section{Constraint set (Brescia 017029, ancoraggio 1/1/2025)}
\label{sec:constraints}
Filosofia: il censimento entra come \emph{distribuzione condizionale applicata ai
conteggi anagrafici} (quote censuarie per gruppo $\times$ totali anagrafici del
gruppo). I marginali demografici restano esatti per costruzione; lo strato
socio-economico eredita la struttura censuaria.

\begin{center}
\small
\begin{tabular}{@{}llp{4.6cm}rr@{}}
\toprule
\textbf{File} & \textbf{Tipo} & \textbf{Contenuto} & \textbf{Celle} & \textbf{Somma}\\
\midrule
c1\_sex\_age\_marital & hard & S$\times$E$\times$stato civile (4 mod.) & 712 & 199\,853\\
c2\_sex\_age\_citizenship & hard$^{*}$ & S$\times$E$\times$ITL/FRG & 395 & 199\,853\\
c3\_sex\_ageclass\_edu & soft & S$\times$classe$\times$istruzione (6 mod.) & 48 & pop $9+$\\
c4\_sex\_ageclass\_condprof & soft & S$\times$classe$\times$condprof (6 mod.) & 48 & 175\,748\\
c5\_edu\_citizenship & soft & ITL/FRG$\times$istruzione & 12 & 186\,538\\
c6\_condprof\_citizenship & soft & ITL/FRG$\times$condprof & 12 & 175\,748\\
nationality\_conditional & two-stage & $P(\text{paese}\mid\text{FRG})$, 143 paesi & --- & ---\\
\bottomrule
\end{tabular}
\end{center}
\noindent{\footnotesize $^{*}$riclassificato hard a valle del risultato in
\S\ref{sec:key-result}.}

\medskip\noindent Collassi applicati (mapping stampati e verificati a runtime):
\begin{itemize}[leftmargin=*]
  \item \textbf{Stato civile} $7\to4$: unioni civili assimilate
  ($15\to$coniugato/unito, $17\to$divorziato/sciolto, $16\to$vedovo).
  \item \textbf{Istruzione} $10\to6$ (label-driven, con ordine dei check):
  nessun\_titolo, elementare, media, diploma (USE\_IF, incl.\ IFTS),
  laurea\_o\_its (BL: ITS + terziario di I livello), post\_laurea (ML\_RDD:
  II livello + dottorato).
  \item \textbf{Condprof} 6 modalit\`a elementari: occupato, in cerca,
  casalinga/o, studente/ssa, percettore pensioni/redditi, altra condizione.
  \item \textbf{Nazionalit\`a}: filtro keyword degli aggregati (``tutte le
  voci'', continenti, UE, \dots); top-5 Brescia: Romania, Pakistan, Ucraina,
  Egitto, India; stranieri $38\,228$ ($19{,}1\%$).
\end{itemize}

%=======================================================================
\section{Risultato chiave: identit\`a anagrafe--censimento}
\label{sec:key-result}
Il report di consistenza confronta cella per cella (sesso$\times$et\`a a singolo
anno) l'anagrafe all'1/1/2025 con il censimento 2024:
\[
\text{MAE} = 0.0, \qquad \text{scarto totale} = 0 \ \text{su } 199\,853.
\]
Le due fonti sono \emph{identiche}: la popolazione post-2018 \`e census-based e
la diffusione SDMX comunale espone la base riconciliata, senza l'arrotondamento
delle tavole web (\S\ref{sec:quirks}, punto 5). Inoltre le stime censuarie sono
internamente additive ($\sigma_{\text{arrotondamento}}=0$ su tutte le tavole con
totali).

\medskip\noindent\textbf{Implicazione metodologica.} La distinzione hard/soft
prevista in fase di design va ridisegnata:
\begin{itemize}[leftmargin=*]
  \item i vincoli demografici (C1, C2) sono di fatto tutti \emph{hard}: stessa
  base dati riconciliata;
  \item il carattere \emph{soft} residuo di C3--C6 non deriva da inconsistenza
  tra fonti (nulla) n\'e da arrotondamenti (nulli), ma esclusivamente dalla
  natura di stima campionaria delle variabili socio-economiche; l'incertezza
  pertinente \`e l'errore campionario documentato da ISTAT per il censimento
  permanente, non misurabile dagli scarti interni;
  \item per il solver: partenza con vincoli hard ovunque, regolarizzazione come
  assicurazione (inversione dell'impostazione iniziale).
\end{itemize}

%=======================================================================
\section{Da fare / prossimi passi}
\begin{itemize}[leftmargin=*]
  \item[$\square$] \textbf{Adattatore GibbsPCD}: da \code{constraints\_2025/} al
  formato vincoli del solver; generazione Brescia v4 (8 attributi, nazionalit\`a
  two-stage).
  \item[$\square$] \textbf{Livello sub-comunale}: fetch \flow{SUB\_MUN\_DATA}
  (tema \code{censpop}) per le 33 aree sub-comunali di Brescia --- sesso, et\`a,
  cittadinanza, istruzione, condizione professionale su basi territoriali 2021;
  attributo zona nel MaxEnt con vincoli della stessa fonte censuaria.
  \item[$\square$] \textbf{$\sigma$ campionarie}: recuperare gli errori di stima
  documentati del censimento permanente per parametrizzare i vincoli soft.
  \item[$\square$] \textbf{Generalizzazione multi-comune}: run su un secondo
  comune (es.\ Torino 001272) per verificare robustezza dei collassi label-driven
  e della gerarchia condprof.
  \item[$\square$] Tavole aggiuntive candidate: famiglie
  (\flow{DF\_DCSS\_FAMIGLIE\_TV\_*}), background migratorio
  (\flow{DF\_DCSS\_MIGR\_BACKG\_*\_COM}), pendolarismo.
\end{itemize}

\end{document}
